% from ega4-18-fr.tex
\oldpage[IV-4]{176}

Or, il résulte de (0, 20.7.15) et de ($0_{\mathrm I}$, 7.2.9) que les $\delta(x)$ engendrent le $k(\mathfrak p)$-espace vectoriel $\hat{\Omega}^1_{A/k}\otimes_A k(\mathfrak p)$; comme par hypothèse $r<m$, il existe donc $z\in A$ tel que $\delta(z)$ ne soit pas combinaison linéaire des $\delta(u_ix_i)$ pour $1\leq i\leq r$. Si $z\notin\mathfrak m$, on prendra $u=z$; sinon, $1+z$ est inversible et on a $\delta(1+z)=\delta(z)$ puisque $d(1+z)=dz$; on prendra alors $u=1+z$, ce qui achève de prouver (18.11.5.2).

Revenons maintenant à la démonstration de l'implication $a)\Rightarrow b)$ dans (18.11.5). Notons d'abord que puisque $\mathfrak p\neq\mathfrak m$, il existe un système de paramètres $(x_i)_{1\leq i\leq n}$ de $A$ (avec $n=\dim(A)$) tel que $x_i\notin\mathfrak p$ pour $1\leq i\leq n$: en effet, on ne peut avoir $x_i\in\mathfrak p$ pour tout $i$, sans quoi $A/\mathfrak p$ serait de longueur finie, et $\mathfrak p$ serait maximal, contrairement à l'hypothèse. Mais si par exemple $x_1\notin\mathfrak p$, il suffit, pour chaque indice $i$ tel que $x_i\in\mathfrak p$, de remplacer $x_i$ par $x_1+x_i$, pour avoir un système de paramètres dont aucun élément n'appartient à $\mathfrak p$. L'hypothèse $a)$ et la relation $\dim(A/\mathfrak q)=n$ entraînent, en vertu de (18.11.5.1), que $m\geq n$; on peut donc considérer une famille $(x_i)_{1\leq i\leq m}$ d'éléments $x_i\notin\mathfrak p$, dont les $n$ premiers forment un système de paramètres de $A$. Multipliant en outre les $x_i$ par des éléments inversibles $u_i$ de $A$, on peut, grâce à (18.11.5.2) supposer que les images des $dx_i$ ($1\leq i\leq m$) dans $(\hat{\Omega}^1_{A/k})_{\mathfrak p}$ engendrent cet $A_{\mathfrak p}$-module, et la multiplication par les $u_i$ n'a pas altéré le fait que les $x_i$ pour $1\leq i\leq n$ forment un système de paramètres.

Considérons alors le $k$-homomorphisme local $u:B\to A$ tel que $u(T_i)=x_i$ pour $1\leq i\leq m$ (Bourbaki, \emph{Alg. comm.}, chap. III, \textsection~4, n\textsuperscript{o}~5, prop. 6); comme les $x_i$ engendrent un idéal de définition de $A$, il résulte de ($0_{\mathrm I}$, 7.4.4 et 7.4.3) que $u$ fait de $A$ une $B$-algèbre finie. On a donc (18.11.3.5) la suite exacte
\[
  \hat{\Omega}^1_{B/k}\otimes_B A\xrightarrow{v}\hat{\Omega}^1_{A/k}\longrightarrow\Omega^1_{A/B}\longrightarrow0.
\]
Mais les $d_Ax_i$ sont les images canoniques par $v$ des éléments $d_BT_i\otimes1$ (0, 20.5.2.6). Si on localise la suite exacte précédente en $\mathfrak p$, on voit donc que $v_{\mathfrak p}:(\hat{\Omega}^1_{B/k}\otimes_B A)_{\mathfrak p}\to(\hat{\Omega}^1_{A/k})_{\mathfrak p}$ est surjectif, et par suite on a $0=(\Omega^1_{A/B})_{\mathfrak p}=\Omega^1_{A_{\mathfrak p}/B_{\mathfrak r}}$ (où $\mathfrak r$ est l'image réciproque de $\mathfrak p$ dans $A$, cf. (16.4.15)). En vertu de (17.4.1), cela entraîne la propriété $b)$ de (18.11.5). C.Q.F.D.

\begin{remark}[18.11.6]
\label{IV.18.11.6}
Dans l'énoncé de (18.11.5), on ne peut supprimer l'hypothèse $\mathfrak p\neq\mathfrak m$. En effet, pour tout $k$-homomorphisme local $u:B\to A$, où $B=k[[T_1,\ldots,T_r]]$ ($r$ entier quelconque) faisant de $A$ une $B$-algèbre finie, $\mathfrak m$ est le seul point de $\Spec(A)$ au-dessus de l'idéal maximal $\mathfrak n$ de $B$; le morphisme $\Spec(A)\to\Spec(B)$ ne peut être non ramifié en $\mathfrak m$ que si $A/\mathfrak nA$ est un corps, extension séparable de $k$ (17.4.1), ce qui implique que le corps résiduel $K$ de $A$ est une extension séparable de $k$. Si cette condition n'est pas remplie, la conclusion de (18.11.5) ne peut jamais être vérifiée pour $\mathfrak p=\mathfrak m$, quel que soit l'entier $m$.
\end{remark}

\begin{corollary}[18.11.7]
\label{IV.18.11.7}
Soient $k$ un corps, $A$ une $k$-algèbre locale noethérienne complète équidimensionnelle, dont le corps résiduel est une extension finie de $k$, $\mathfrak p$ un idéal premier de $A$. Alors le nombre minimum de générateurs du $A_{\mathfrak p}$-module $(\hat{\Omega}^1_{A/k})_{\mathfrak p}$ est au moins égal à $\dim(A)$.
\end{corollary}

C'est un cas particulier de (18.11.5.1).

Plus particulièrement:

